\documentclass[12pt]{letter}

\usepackage[margin=1.25in]{geometry}
\usepackage[hidelinks]{hyperref}
\usepackage{setspace}

\emergencystretch=3em

\signature{Mathias Valla\\Chaire DIALog, Institut Louis Bachelier\\\href{mailto:mathias.valla@gmail.com}{mathias.valla@gmail.com}}
\address{Mathias Valla\\Chaire DIALog\\Institut Louis Bachelier\\Paris, France\\\href{mailto:mathias.valla@gmail.com}{mathias.valla@gmail.com}}

\begin{document}

\begin{letter}{Editors\\Journal of Classification}

\opening{Dear Editors,}

Please consider the manuscript ``Spherical Splits for Decision Trees and
Random Forests: A Short Note'' for publication in the Journal of
Classification as a short note.

The manuscript proposes and evaluates a simple geometric extension of
tree-based classification and regression models: replacing some
axis-aligned threshold questions by hyperspherical membership questions. I
believe the note fits the Journal of Classification because it concerns
classification methodology, tree-structured models, statistical learning, and
the geometry of recursive partitions. The contribution is intentionally
compact: it defines the split rule, records several useful theoretical
properties, notes that CART cost-complexity pruning transfers directly, and
reports exploratory evidence that spherical random forests can improve mean
classification performance on tabular benchmark datasets.

The manuscript is original work, has not been published previously, and is
not under consideration elsewhere. To respect the journal's double-blind
review policy, the main manuscript is anonymized. Author details and
declarations are supplied on a separate title page. A blinded reproduction
document and supplementary code/results are prepared for review; the public
repository is \url{https://github.com/MathiasValla/spherical-decision-trees}
and can be disclosed in the published version.

\closing{Sincerely,}

\end{letter}

\end{document}
